\chapter{Algorithmes : calculer, prouver et contrôler l’erreur}
\section*{Exercice 24.1 -- Invariant de boucle}
[Programmation] Écrire une boucle calculant $n!$ pour $n\ge0$ et justifier sa correction et sa terminaison.

\section*{Exercice 24.2 -- Coût de la dichotomie}
[Méthode] Sur un intervalle de longueur 3, combien de divisions suffisent pour une erreur du milieu au plus $10^{-6}$ ?

\section*{Exercice 24.3 -- Newton cyclique}
[Critique] Pour $f(x)=x^3-2x+2$, calculer deux itérations de Newton depuis 0. Quelle protection prévoir dans un programme ?

\section*{Exercice 24.4 -- Erreur des rectangles}
[Preuve] Pour $\ln2=\int_1^2dx/x$, donner un encadrement par $n$ rectangles et la largeur de cet encadrement.

\section*{Exercice 24.5 -- Monte-Carlo et erreur}
[Application] On estime $\pi$ par $\widehat\pi_n=4S_n/n$, où $S_n$ compte les points uniformes indépendants tombant dans le quart de disque unité. Majorer la probabilité d'une erreur au moins $\varepsilon$.

\section*{Exercice 24.6 -- Résidu et erreur}
[Raisonnement] Si $f(x)=10^{-9}(x-1000)$, un résidu $|f(0)|=10^{-6}$ signifie-t-il que 0 est proche de la racine ?
